\documentclass[10pt,a4paper]{article}
\usepackage[T1]{fontenc}
\usepackage{lmodern}
\usepackage[margin=23mm]{geometry}
\usepackage{amsmath,amssymb,booktabs,microtype}
\usepackage[hidelinks]{hyperref}
\pdfinfoomitdate=1
\pdftrailerid{}
\pdfsuppressptexinfo=15
\setlength{\parindent}{0pt}
\setlength{\parskip}{5pt}
\title{Prescribed-temperature Parker winds\\with an optional linear stellar tide}
\author{\texttt{exojax-escape}: model definition and verification}
\date{Definition revision 1; 21 November 2026}
\begin{document}
\maketitle

\section{What this model calculates}
Given a temperature profile $T(r)$, a planet mass and a constant mean molecular
weight $\mu$, the model finds a smooth radial wind that passes from below to above
the local sound speed. It can add the leading stellar tidal acceleration. It does
\emph{not} calculate the temperature from heating and cooling, or update $\mu$
from ionization. These are explicit physical approximations, even when the
numerical equations are solved accurately.

\[
\boxed{T(r),\ \mu,\ M_{\rm p},\ \kappa}
\ \longrightarrow\ \boxed{\text{critical radius and regular wind}}
\ \longrightarrow\ \boxed{v(r)\ \text{and parameter derivatives}}
\]

The existing \texttt{parker.py} solves the constant-temperature, no-tide case.
The existing \texttt{nonisothermal.py} solves a supplied smooth temperature
profile, with or without tides. This document records both implementations
at source revision \texttt{310a434}; it introduces no new wind physics.

\section{Pressure, gravity and the tidal approximation}
Write the ideal-gas pressure as $P=\rho a^2$, where
$a^2=k_{\rm B}T/(\mu m_{\rm p})$. Here $a$ is the local \emph{isothermal}
sound speed; no adiabatic index is inserted. Steady spherical continuity gives
$4\pi r^2\rho v=\dot M$. Eliminating the density gradient from radial momentum
balance gives
\begin{equation}
 \left(v-\frac{a^2}{v}\right)\frac{dv}{dr}
 =\frac{2a^2}{r}-\frac{da^2}{dr}-\frac{GM_{\rm p}}{r^2}+\kappa r .
 \label{eq:explanation-momentum}
\end{equation}
The term $-da^2/dr$ is necessary when temperature changes with radius. The
pressure/continuity reduction agrees with the standard isothermal development
in Cranmer~\cite{cranmer}; retaining this derivative is our explicit extension.
The velocity equation does not fix the density normalization or mass-loss rate.

In the leading expansion of the rotating stellar potential along the
star--planet axis, for $r\ll a_{\rm orb}$ and stellar mass much greater than planet
mass, the differential stellar gravity contributes $2GM_\star r/a_{\rm orb}^3$
and the centrifugal term contributes $GM_\star r/a_{\rm orb}^3$. Thus
\begin{equation}
 \kappa=\frac{3GM_\star}{a_{\rm orb}^3},\qquad
 g_{\rm eff}=\frac{GM_{\rm p}}{r^2}-\kappa r .
\end{equation}
The tide reduces inward gravity. The API takes \texttt{tidal\_coefficient\_s2}
as an input; it does not calculate the orbit. The official p-winds isothermal
tidal implementation uses this same coefficient~\cite{pwinds}. We inspected
its equations but did not execute p-winds as a benchmark. Retaining spherical
continuity while using an axis-dependent tidal term is a reduced model,
\emph{not} a full multidimensional Roche flow.

\section{Crossing the sonic point without a singularity}
At $v=a$, the factor multiplying $dv/dr$ in
Eq.~\eqref{eq:explanation-momentum} vanishes. A regular wind must make its
right-hand side vanish at the same radius. The solver first locates this
critical point; it then enforces the finite accelerating slope there.

The logarithmic variables $x=\ln(r/R)$ and $q=\ln(v/a)$ give
\begin{equation}
 (e^{2q}-1)(q_x+\beta/2)=N(x),\qquad
 \beta=\frac{d\ln a^2}{dx},\quad
 N=2-\beta-\frac{GM_{\rm p}}{ra^2}+\frac{\kappa r^2}{a^2}.
\end{equation}
At the critical point $q=N=0$. Differentiating the equation there yields
$2q_x^2+\beta q_x-N_x=0$; the selected positive root is recorded in the
appendix. The velocity slope is $d\ln v/d\ln r=q_x+\beta/2$, so it is not in
general equal to the Mach-number slope.

The computation starts from the isothermal no-tide wind at $T(r_{\rm out})$.
Eight continuation steps introduce the temperature variation and tide. Two
Chebyshev grids meet at the critical point and carry one continuous wind.
The final equation is solved by damped Newton iteration. Both the desired
critical point and the isothermal initializer critical point must lie inside
the supplied interval; this is a limitation of this solver, not a theorem
about all possible winds. Invalid or unconverged solutions return NaNs.

\section{What is differentiated and what has been verified}
The sonic radius is itself parameter dependent. A change in temperature or
tide moves that radius and both numerical grids. The implementation uses two
\texttt{jax.lax.custom\_root} calls: a scalar critical-point solve and a vector
wind solve. Their implicit derivatives include this moving grid and the final
interpolation. Differentiation applies to the converged discrete equations;
it requires an isolated root and a nonsingular residual Jacobian.

All pressure and gravity calculations use cgs units. Both public velocity
functions return \textbf{km/s}, dividing the internal cm/s velocity by $10^5$.
The equation bindings record this conversion separately, so the returned
quantity cannot be mistaken for the internal velocity.

\begin{center}
\begin{tabular}{lll}
\toprule
Question & MCEG assessment & Basis \\
\midrule
Do the equations describe the code? & supported & Algebra/source review \\
Do the recorded numerical checks pass? & supported & 25 retained test cases \\
Is the atmospheric approximation accurate? & unknown & No matched validation \\
\bottomrule
\end{tabular}
\end{center}

The tests compare against independent Lambert-$W$ evaluation, independent
SciPy sonic shooting, a constant-temperature tidal integral and a cooling
power-law critical radius. They also check sonic regularity, derivatives
against finite differences, forward/reverse differentiation, moving-sonic-point
continuity, and 129/257-node refinement. The exact inputs and tolerances remain
in the pinned test files; these sampled checks do not define a continuous
accuracy envelope. The independent shooting calculation solves the same
physical ODE, so agreement verifies the numerics, not the atmospheric closure.
No Cloudy wind comparison or NUTS posterior validation is claimed.

\texttt{mceg/contracts/escape.parker\_wind.yaml} is the equation record.
Its bindings point to named functions and SHA-256-pinned source files.
This PDF includes the generated equation appendix rather than maintaining a
second mathematical definition. The compact execution record is in
\path{mceg/evidence/parker_wind_20261121/}: command/environment/source hashes,
test log, JUnit result and review. MCEG checks the records and their identity;
it does not prove the equations or replace the scientific review. Runtime
\texttt{escape} code does not import or depend on MCEG.

\begin{thebibliography}{9}
\bibitem{cranmer} S. R. Cranmer (2004), ``New views of the solar wind with the
Lambert W function,'' \emph{American Journal of Physics} 72, 1397.
\href{https://doi.org/10.1119/1.1775242}{doi:10.1119/1.1775242};
\href{https://arxiv.org/pdf/astro-ph/0406176}{author manuscript}, Eqs.~9, 11--16.
\bibitem{pwinds} p-winds contributors, official
\href{https://p-winds.readthedocs.io/en/latest/_modules/parker.html}{\texttt{parker} source}:
\texttt{radius\_sonic\_point\_tidal} and \texttt{structure\_tidal},
inspected 21 November 2026. Implementation reference, not an executed benchmark.
\end{thebibliography}

\clearpage
\appendix
\section{Recorded equations and units}
\small
The following definitions are generated from the MCEG contract. The source
code bindings and assessment pins remain in that contract.
% Generated by mceg equations; edit the contract, then regenerate.
% Requires amsmath. Authored math is retained; TeX is not executed here.
% Definitions and hashes do not establish scientific validity.

% equation-sha256: 5f1b83db451a1beddd5bebac5695edc4e41446ade8bd628210b6d0fbbab65238
\paragraph{Equation wind\_momentum (revision 1)}
Pressure, steady spherical continuity and prescribed-temperature momentum balance.
\begin{equation}
\begin{aligned}
a^2=\frac{k_{\rm B}T(r)}{\mu m_{\rm p}},\quad P=\rho a^2,\quad 4\pi r^2\rho v=\dot M,\\
\kappa=\frac{3GM_\star}{a_{\rm orb}^3},\quad g_{\rm eff}=\frac{GM_{\rm p}}{r^2}-\kappa r,\\
\left(v-\frac{a^2}{v}\right)\frac{dv}{dr}=\frac{2a^2}{r}-\frac{da^2}{dr}-g_{\rm eff}
\end{aligned}
\label{eq:escape.parker_5f_wind:wind_5f_momentum:1}
\end{equation}
\noindent\textbf{Scope:} Constant mu; optional linear tide along the star-planet axis, combined with spherical continuity.

\begin{description}
\item[$r$] Radius (cm; ID: radius).
\item[$T$] Prescribed temperature (K; ID: temperature).
\item[$\mu$] Mean particle mass in proton-mass units (1; ID: mu).
\item[$M_{\rm p}$] Planet mass (g; ID: mass).
\item[$\kappa$] Supplied tidal coefficient (s\textasciicircum{}-2; ID: tide).
\item[$v$] Internal radial velocity (cm s\textasciicircum{}-1; ID: velocity).
\item[$a$] Local isothermal sound speed (cm s\textasciicircum{}-1; ID: sound\_speed).
\item[$\rho$] Mass density; normalization not determined by the velocity solver (g cm\textasciicircum{}-3; ID: density).
\end{description}

\noindent\textbf{Assumptions:}
\begin{itemize}
\item The caller supplies kappa; the wind API does not infer stellar mass or orbital separation.
\item The stellar-mass-dominated, r much smaller than orbital separation expansion is not full Roche geometry.
\item Temperature is supplied and mu is fixed. No energy or ionization closure is solved.
\end{itemize}

% equation-sha256: 9d848f938218e7c751318220f2305fc9c64d21b847ff2a201a66de2b90d3dcca
\paragraph{Equation isothermal\_limit (revision 1)}
No-tide isothermal Parker branch and its regular derivative.
\begin{equation}
\begin{aligned}
r_{\rm s}=\frac{GM_{\rm p}\mu m_{\rm p}}{2k_{\rm B}T},\quad t=\ln(r/r_{\rm s}),\quad z=\ln(v/a),\\
e^{2z}-2z-1=4(t+e^{-t}-1),\quad \operatorname{sign}z=\operatorname{sign}t,\\
\frac{dz}{dt}=\frac{2(1-e^{-t})}{e^{2z}-1},\qquad \left.\frac{dz}{dt}\right|_{t=0}=1
\end{aligned}
\label{eq:escape.parker_5f_wind:isothermal_5f_limit:1}
\end{equation}
\noindent\textbf{Scope:} The analytic-limit implementation solves the scalar implicit relation by bisection; it does not call Lambert W.

\begin{description}
\item[$r$] Radius (cm; ID: radius).
\item[$r_{\rm s}$] Isothermal sonic radius (cm; ID: sonic\_radius).
\item[$T$] Constant temperature (K; ID: temperature).
\item[$M_{\rm p}$] Planet mass (g; ID: mass).
\item[$\mu$] Mean particle mass in proton-mass units (1; ID: mu).
\item[$t$] Log radius divided by sonic radius (1; ID: coordinate).
\item[$z$] Log Mach number (1; ID: log\_mach).
\end{description}

\noindent\textbf{Assumptions:}
\begin{itemize}
\item Float64 and positive inputs; the accelerating wind branch is selected.
\item A fifth-order sonic series for |t|<1e-3 avoids cancellation; custom JVP uses its derivative.
\end{itemize}

% equation-sha256: b9620ae024389487663d0f4c42b673fa718090a82bb0692eb89577063e07a8f3
\paragraph{Equation dimensionless\_wind (revision 1)}
Dimensionless residual and continuation from the isothermal no-tide wind.
\begin{equation}
\begin{aligned}
x=\ln(r/R),\quad L=\ln(r_{\rm out}/R),\quad a_{\rm ref}^2=\frac{k_{\rm B}T(r_{\rm out})}{\mu m_{\rm p}},\\
s_\lambda=1+\lambda\left[\frac{T(Re^x)}{T(r_{\rm out})}-1\right],\quad \beta_\lambda=\frac{\partial\ln s_\lambda}{\partial x},\\
A=\frac{GM_{\rm p}}{Ra_{\rm ref}^2},\quad B=\frac{\kappa R^2}{a_{\rm ref}^2},\quad q=\ln\frac{v}{a_{\rm ref}\sqrt{s_\lambda}},\\
N_\lambda=2-\beta_\lambda-\frac{Ae^{-x}}{s_\lambda}+\frac{\lambda Be^{2x}}{s_\lambda},\\
(e^{2q}-1)\left(q_x+\frac{\beta_\lambda}{2}\right)-N_\lambda=0
\end{aligned}
\label{eq:escape.parker_5f_wind:dimensionless_5f_wind:1}
\end{equation}
\noindent\textbf{Scope:} The target model is lambda=1. Intermediate lambda values are numerical initialization only.

\begin{description}
\item[$r$] Radius (cm; ID: radius).
\item[$R$] Inner domain radius (cm; ID: base).
\item[$r_{\rm out}$] Outer domain radius (cm; ID: outer).
\item[$x$] Log radius relative to the inner boundary (1; ID: coordinate).
\item[$\lambda$] Continuation fraction (1; ID: continuation).
\item[$s_\lambda$] Squared sound speed divided by its reference value (1; ID: sound\_ratio).
\item[$q$] Log Mach number (1; ID: log\_mach).
\item[$N_\lambda$] Dimensionless momentum numerator (1; ID: numerator).
\end{description}

\noindent\textbf{Assumptions:}
\begin{itemize}
\item Eight continuation steps; default 129 total collocation nodes.
\item Reference temperature is evaluated at the outer boundary, not at the sonic point.
\end{itemize}

% equation-sha256: 722d8907a85c6c31aa007247f25cd41aa148bbbda06f04bb3dfdb7fe39c7cd4f
\paragraph{Equation sonic\_regularity (revision 1)}
Regular transonic critical point and slope.
\begin{equation}
\begin{aligned}
N_\lambda(x_{\rm c})=0,\quad q(x_{\rm c})=0,\quad N_{\lambda,x}(x_{\rm c})>0,\\
q_{x,{\rm c}}=\frac{-\beta_{\lambda,{\rm c}}+\sqrt{\beta_{\lambda,{\rm c}}^2+8N_{\lambda,x}(x_{\rm c})}}{4},\\
\left.\frac{d\ln v}{d\ln r}\right|_{\rm c}=q_{x,{\rm c}}+\frac{\beta_{\lambda,{\rm c}}}{2}
\end{aligned}
\label{eq:escape.parker_5f_wind:sonic_5f_regularity:1}
\end{equation}
\noindent\textbf{Scope:} One negative-to-positive numerator root in the supplied radial interval.

\begin{description}
\item[$x_{\rm c}$] Log critical radius (1; ID: critical).
\item[$N_\lambda$] Dimensionless numerator (1; ID: numerator).
\item[$\beta_\lambda$] Log sound-speed-squared slope (1; ID: beta).
\item[$q$] Log Mach number (1; ID: log\_mach).
\end{description}

\noindent\textbf{Assumptions:}
\begin{itemize}
\item The implementation also requires subsonic inner and supersonic outer branches.
\item The no-tide isothermal initializer critical point must lie inside the domain.
\end{itemize}

% equation-sha256: 3c88e12900ef54fcc8caeed05c93ea579519ec2b18a330d6ed45a7cfc1098ff6
\paragraph{Equation collocation (revision 1)}
Two Chebyshev-Lobatto pieces joined at the critical point.
\begin{equation}
\begin{aligned}
\xi_i=\frac{1-\cos(\pi i/(m-1))}{2},\quad i=0,\ldots,m-1,\\
x_i^-=x_{\rm c}\xi_i,\qquad x_i^+=x_{\rm c}+(L-x_{\rm c})\xi_i,\\
\mathcal R_i=(e^{2q_i}-1)\left[(D_xq)_i+\frac{\beta_i}{2}\right]-N_i=0,\\
\mathcal R_{\rm c}\leftarrow q_{\rm c},\quad
\mathcal R_0\leftarrow(D_x^-q)_{\rm c}-q_{x,{\rm c}},\quad
\mathcal R_{n-1}\leftarrow(D_x^+q)_{\rm c}-q_{x,{\rm c}}
\end{aligned}
\label{eq:escape.parker_5f_wind:collocation:1}
\end{equation}
\noindent\textbf{Scope:} The repeated critical node is stored once: total n=2m-1. The displayed first/last replacements are residual rows, not imposed endpoint velocities.

\begin{description}
\item[$n$] Odd total node count; at least 17 (1; ID: grid\_size).
\item[$x$] Log radial coordinate (1; ID: coordinate).
\item[$x_{\rm c}$] Log critical radius (1; ID: critical).
\item[$D_x$] Barycentric derivative matrix rescaled to each log-radius interval (1; ID: derivative).
\item[$q$] Collocated log Mach number (1; ID: log\_mach).
\item[$\mathcal R$] Discrete momentum and regularity residual (1; ID: residual).
\end{description}

\noindent\textbf{Assumptions:}
\begin{itemize}
\item A damped Newton solve follows continuation; accepted residual infinity norm must be below 1e-8.
\item This is a discretized ODE, not an exact continuum solution.
\end{itemize}

% equation-sha256: d9d27bdf5b0169c836c035a3c92462862d78924cac116031f465ebde2a82fce6
\paragraph{Equation implicit\_derivative (revision 1)}
Implicit derivatives of the critical point and the converged discrete wind.
\begin{equation}
\begin{aligned}
\frac{dx_{\rm c}}{d\theta}=-\frac{N_\theta}{N_x},\qquad
\frac{dq}{d\theta}=-\mathcal R_q^{-1}\left(\mathcal R_\theta+\mathcal R_{x_{\rm c}}\frac{dx_{\rm c}}{d\theta}\right)
\end{aligned}
\label{eq:escape.parker_5f_wind:implicit_5f_derivative:1}
\end{equation}
\noindent\textbf{Scope:} At lambda=1; partial derivatives include prescribed-temperature parameters, mass, mu, tide and domain bounds.

\begin{description}
\item[$\theta$] One physical input parameter; units depend on parameter (parameter dependent; ID: parameters).
\item[$x_{\rm c}$] Log critical radius (1; ID: critical).
\item[$q$] Vector of collocated log Mach numbers (1; ID: log\_mach).
\item[$\mathcal R$] Discrete residual including both moving grids (1; ID: residual).
\end{description}

\noindent\textbf{Assumptions:}
\begin{itemize}
\item jax.lax.custom\_root differentiates the residual equations, not bisection or Newton iterations.
\item The scalar critical slope and the discrete residual Jacobian must be nonsingular.
\item The output interpolation and sound-speed factor are differentiated as well.
\end{itemize}

% equation-sha256: 897616fa9ca51dd9bc6361dede83dc37bf52202d037e3b175b51486a3f95db10
\paragraph{Equation velocity\_output (revision 1)}
Public output converts cgs sound speed and velocity to km/s.
\begin{equation}
\begin{aligned}
v_{\rm out}(r)=10^{-5}\sqrt{\frac{k_{\rm B}T(r)}{\mu m_{\rm p}}}\,\exp q(\ln(r/R))\quad[\mathrm{km\,s^{-1}}]
\end{aligned}
\label{eq:escape.parker_5f_wind:velocity_5f_output:1}
\end{equation}
\noindent\textbf{Scope:} For the analytic isothermal function, use z(ln(r/r\_s)) in place of q(ln(r/R)), with constant T.

\begin{description}
\item[$r$] Requested radius (cm; ID: radius).
\item[$T$] Prescribed or constant temperature (K; ID: temperature).
\item[$\mu$] Mean particle mass in proton-mass units (1; ID: mu).
\item[$v_{\rm out}$] Returned radial velocity (km s\textasciicircum{}-1; ID: velocity\_kms).
\item[$q$] Interpolated log Mach number (1; ID: log\_mach).
\end{description}

\noindent\textbf{Assumptions:}
\begin{itemize}
\item Nonisothermal values use piecewise Chebyshev interpolation inside the domain.
\item Invalid or unconverged nonisothermal solutions return NaN; diagnostics are available with return\_info=True.
\end{itemize}

\end{document}
